\section{Euler localization and support collisions}
\label{sec:elliptic-hall-localization}

Direct sum also gives a different coproduct after multiplication by
the Euler class of the curve Hom complex. The global diagonal
determines this factor. It produces coherent identities in explicit
localizations, but it changes the coproduct even on the constant
charge-two class.

\subsection{The global Hom complex}

On $E_1\times E_2\times B\mathbb G_m^2$, let $L_i$ be the two
universal line bundles and let $D_i$ be the graph of the $i$th point
in the additional geometric curve. The universal length-one sheaf
is $\mathcal O_{D_i}\otimes L_i$.
Put $L=L_1^\vee\otimes L_2$ and $w=z_2-z_1=c_1(L)$.
Resolve the first sheaf by
$[\mathcal O(-D_1)\otimes L_1\to\mathcal O\otimes L_1]$ in
degrees $-1,0$. Applying Hom into the second and pushing along the
geometric curve gives the actual perfect complex
\begin{equation}\label{gh:global-hom}
 [\,L\longrightarrow L(\Delta_{12})\,]_{0,1}.
\end{equation}
Its differential is the canonical diagonal section. It is acyclic
away from the diagonal and has both Hom and Ext$^1$ at coincident
points. Thus replacing it by its generic value zero would lose the
collision.

Its K-class is $[L]-[L(\Delta_{12})]$, so its localized Euler class
and Chern ratio are
\begin{equation}\label{gh:euler}
 e(D_{12})=\frac{w}{w+\delta_{12}}
        =1-\frac{\delta_{12}}w,\qquad
 \Psi(D_{12},t)=\frac{t+w}{t+w+\delta_{12}}
        =1-\frac{\delta_{12}}{t+w}.
\end{equation}
Here $\delta_{12}^2=\Delta_*(c_1(T_E))=0$.
The odd cohomology terms of the diagonal in
\eqref{gh:diagonal} are part of these formulas.

\subsection{An Euler-corrected coproduct}

For ordered blocks $A,B$ define
\[
 D_{A,B}=\prod_{i\in A,\ j\in B}(z_j-z_i),\qquad
 E_{A,B}=K_{A,B}^{-1}.
\]
The cross resultant $D_{A,B}$ is invariant under permutations
inside each block. It is a nonzero divisor on the polynomial
module with exterior coefficients. Inverting it also inverts
each $z_j-z_i+\delta_{ij}$, by a finite nilpotent expansion.
The universal curve Hom complex on
$\mathcal M_a\times\mathcal M_b$ restricts under ordered direct
sum to the sum of the complexes \eqref{gh:global-hom}.
Thus $E_{A,B}$ is its actual localized Euler class.

Define
\begin{equation}\label{gh:localized-coproduct}
 \Delta^{\mathrm{loc}}_{a,b}(F)
   =E_{A,B}\,s_{a,b}^*F
   \ \in\
   (\mathsf B_a\otimes\mathsf B_b)[D_{A,B}^{-1}].
\end{equation}
This uses ordinary cohomological direct-sum pullback and the
specified Euler factor. It has no additional exceptional trace.
For several blocks, comparisons are made in the common
localization by all cross resultants that occur.

Let $\sigma_{A,B}$ be the ordinary Koszul block swap, including
the Chern variables. The localized crossing is
\begin{equation}\label{gh:localized-crossing}
 \beta^{\mathrm{loc}}_{A,B}
    =\sigma_{A,B}\circ M_{K_{A,B}/K_{B,A}},
\end{equation}
where $M_f$ denotes multiplication by $f$.
Equivalently, with the swapped output labels fixed, it is
$M_{K_{B,A}/K_{A,B}}\circ\sigma_{A,B}$.
For one label in each block this output formula is
$(1-2\delta_{12}/w)\sigma$.
The side on which the multiplier acts is part of the formula.

\begin{proposition}\label{gh:localized-identities}
The maps \eqref{gh:localized-coproduct} are coassociative with
identity zero-endpoint components and charge-zero counit.
With the crossings \eqref{gh:localized-crossing}, they satisfy
the bialgebra interchange identities for the actual Hall product
in the stated common localizations, in every charge.
The localized crossings satisfy the block hexagons, braid
identities and involution.
\end{proposition}
\begin{proof}
For three blocks,
\[
 E_{A,B}E_{A\cup B,C}
 =E_{A,B}E_{A,C}E_{B,C}
 =E_{B,C}E_{A,B\cup C}.
\]
Together with direct-sum associativity, this proves coassociativity.
An empty block has Euler factor one and its direct-sum map is the
identity.

The ratios of $K$ factors multiply under disjoint block unions.
This proves the two hexagons. On three blocks both braid
composites have the same permuted, even scalar factors.
Swapping the two blocks exchanges numerator and denominator,
proving the involution.

For interchange, split the first input into blocks $A,B$ and the
second into blocks $C,D$. The outputs receive $A\cup C$ and
$B\cup D$. A term of
$\Delta^{\mathrm{loc}}(F*G)$ has coefficient
\[
 \begin{split}
 \frac{K_{A\cup B,C\cup D}}{K_{A\cup C,B\cup D}}
 &=\frac{K_{A,C}K_{A,D}K_{B,C}K_{B,D}}
         {K_{A,B}K_{A,D}K_{C,B}K_{C,D}}\\
 &=\frac{K_{A,C}K_{B,D}}{K_{A,B}K_{C,D}}\,
                           \frac{K_{B,C}}{K_{C,B}}.
 \end{split}
\]
On the other route, the two coproducts contribute
$1/(K_{A,B}K_{C,D})$, the middle crossing contributes
$K_{B,C}/K_{C,B}$ in these labels, and the output Hall products
contribute $K_{A,C}K_{B,D}$.
The coefficients agree.
Choosing the two input splits and then the two output shuffles
is the same indexing as an input shuffle followed by the output
split. The resulting permutations, and therefore all Koszul
signs, agree.

Initially make the comparison after inverting every Chern
difference on these labels. The left side lies in the stated
output localization, so the identity shows cancellation of the
additional intermediate denominators in the complete sum.
Localization is injective on the polynomial ring tensored with
$H^*(E^n;\mathbb Q)$, since it is a free module over the
polynomial ring. Nilpotent or odd cohomological coefficients do
not change this argument. This proves the asserted identities
without claiming that each intermediate term is polynomial.
\end{proof}

The proposition specifies maps and their identities in localized
tensor recipients. It does not assert an antipode in an unstated
completion. Its crossing is a rational multiplier followed by a
swap, whereas \eqref{gh:polynomial-R} is a polynomial
divided-difference operator.

\begin{example}[Two coproducts on the same Hall algebra]
\label{gh:two-coproducts}
The geometric direct-sum coproduct of
Theorem~\ref{gh:global-hopf} and the Euler-corrected coproduct
satisfy
\[
 d_{1,1}(\mathbf1_2)=1\otimes1,\qquad
 \Delta^{\mathrm{loc}}_{1,1}(\mathbf1_2)
                   =1-\frac{\delta_{12}}{z_2-z_1}.
\]
Thus the identity on the Hall carrier is not a coalgebra
comparison between these maps into the same localized tensor.
No assertion about a different tensor functor follows from this
calculation.
The relation $r=z_1-z_2+\delta_{12}$ of
Proposition~\ref{gh:ordinary-obstruction} obeys
$\beta^{\mathrm{loc}}r=-r$; hence
$(1+\beta^{\mathrm{loc}})r=0$.
This explains how the localized crossing avoids that ordinary
interchange obstruction.

The degree-zero quotient of Corollary~\ref{gh:fibre-quotient}
does not extend coefficient-linearly through this tensor
localization: it sends $z_2-z_1$ to zero, while that element has
been made invertible. The ordinary polynomial quotient belongs
to the polynomial-crossing construction with the direct-sum
coproduct.
\end{example}

\subsection{Nonsplit fibres and support parameters}

The global operation comparison includes the proper flag traces,
but it is a statement after direct image to $BE$. It does not by
itself identify an operation over a moving support divisor.
The following local models display the coefficient and trace
information that such a comparison must retain.

After fixing the first point by translation, let
$\delta=e_2-e_1\in E$. The position part of
\eqref{gh:global-hom} is
\[
 [\,\mathcal O_E\xrightarrow{s_0}\mathcal O_E(0)\,]_{0,1},
\]
with $s_0$ vanishing at the origin. Its derived fibre has one copy
of $\mathbb C$ in degrees zero and one at the collision and is
acyclic elsewhere. Its determinant is $\mathcal O_E(-0)$.
The $S$-factor contributes the even degree-two copy, so the
threefold cross determinant has degree $-2$.
For a split pair the two opposite cross blocks give determinant
degree $-4$ and root degree $-2$.
Ordinary restriction of the degree-one skyscraper representing
this complex would omit its degree-zero Tor term.

On a local parameter chart the extension family is
\[
 Z=\begin{pmatrix}y&t\\0&x\end{pmatrix},\qquad
 g=\begin{pmatrix}1&h\\0&1\end{pmatrix},\qquad
 g^{-1}Zg=
 \begin{pmatrix}y&t+(y-x)h\\0&x\end{pmatrix}.
\]
With endpoint trivializations fixed, the extension groupoid is
$[\mathbb A^1_t/\mathbb G_a]$, where
$h$ acts by $t\mapsto t+(y-x)h$.
Away from $x=y$ it is a point; at $x=y$ it is
$\mathbb A^1\times B\mathbb G_a$.
The nonzero values of $t$ there are nonsplit extensions.
Its dimension is zero throughout, as required by
Proposition~\ref{gh:sector}.

\begin{example}[The binary Quot fibre]\label{gh:double-fibre}
For $T=\mathbb C[\zeta]/\zeta^2$, choose
$Z=\left(\begin{smallmatrix}0&1\\0&0\end{smallmatrix}\right)$.
A line $[u:v]$ is $Z$-invariant precisely when
\[
 \det\begin{pmatrix}u&v\\v&0\end{pmatrix}=-v^2=0.
\]
The actual binary Quot fibre is therefore the double point
$\operatorname{Spec}\mathbb C[v]/(v^2)$ in the chart $u=1$.
It has one geometric point and length two.
For two distinct support points the fibre consists of two
reduced points, whereas for $k_e\oplus k_e$ it is
$\mathbb P^1$.
These are fibres of the same proper correspondence.
The coefficient map is the identity, while its proper trace on
the constant class is multiplication by two, by
Proposition~\ref{gh:constants}. The trace is not the count of
geometric points in a singular fibre.
\end{example}

The full flag coefficient map also exists on a nonsplit
three-step family. On $T=\mathbb A^3_{r,s,t}$ take
\[
 Z=\begin{pmatrix}0&r&t\\0&0&s\\0&0&0\end{pmatrix}.
\]
Since $Z^3=0$, this is a flat rank-three family on the third
infinitesimal neighbourhood of an actual point of $E$.
The first line and first plane are invariant and give a full
flag. At $(r,s,t)=(1,1,0)$ it is the cyclic module
$\mathbb C[\zeta]/\zeta^3$, while at the origin it is split.
The entry $(Z^2)_{13}=rs$ distinguishes them.
Proposition~\ref{gh:coefficient-maps} gives two composites equal
to $\operatorname{id}_{\mathbb Q_T}$ on this family, with their
specified common source and recipient. Changing the sign of
only the $(2,1)$ coefficient map changes one composite by
$-2\operatorname{id}_{\mathbb Q_T}$.
Thus a coefficient flag equation has content beyond its binary
and unit normalizations.

The proper trace of three constant inputs is
$6\mathbf1_3$, and the global four-block identity is proved in
Theorem~\ref{gh:global-hopf}. To compare the operations over moving
support, the scalar trace must be replaced by its coefficient
sheaf and its collision maps. Section~\ref{sec:companion-family}
computes these maps on the cyclic length-three family and proves
an obstruction to extending the proper trace to its ordinary
decomposition image. Section~\ref{sec:support-coalgebra} then
constructs the coefficient coalgebra over support.
Section~\ref{sec:supported-springer} supplies the supported
crossing and Hall interchange from proper flags and compares
their actual evaluations on that family. The strict-splitting
trace obstruction remains in its stated recipient.
